% Figure: an acid-base titration curve.
% pH (a logarithmic quantity) against volume of added base for a weak
% acid titrated by a strong base. The curve is sigmoidal, with a flat
% buffer region where pH = pKa at the half-equivalence point, and a steep
% jump at the equivalence point. Drawn with an explicit logistic-like
% shape (no chemistry solver; a smooth analytic stand-in for the curve).
\begin{figure}[htbp]
\centering
\resizebox{0.82\linewidth}{!}{%
\begin{tikzpicture}
\begin{axis}[
  width=12cm, height=7.4cm,
  xlabel={volume of base added $V$ (mL)},
  ylabel={pH $= -\log_{10}[\mathrm{H}^{+}]$},
  xmin=0, xmax=50, ymin=2, ymax=13,
  xtick={0, 10, 20, 25, 30, 40, 50},
  ytick={2, 4, 6, 8, 10, 12},
  axis lines=left,
  tick align=outside,
  every axis plot/.append style={very thick},
  label style={font=\footnotesize},
  tick label style={font=\footnotesize},
  clip=false,
]
  % Analytic stand-in for a weak-acid/strong-base curve:
  % pH(V) = 3 + 9 / (1 + exp(-0.55*(V - 25))) plus a gentle buffer tilt.
  % Equivalence point at V = 25 mL (steepest region), pKa read at V = 12.5.
  \addplot[engblue, domain=0:50, samples=250]
    {3 + 9/(1 + exp(-0.55*(x - 25))) + 0.02*(x-25)};
  % half-equivalence marker: V = 12.5 mL, pH = pKa
  \addplot[engred, only marks, mark=*, mark size=1.6pt]
    coordinates {(12.5, 4.75)};
  \draw[engred, dotted, thick] (axis cs:12.5, 2) -- (axis cs:12.5, 4.75);
  \draw[engred, dotted, thick] (axis cs:0, 4.75) -- (axis cs:12.5, 4.75);
  \node[font=\footnotesize, align=left, text=engred, anchor=west]
    at (axis cs:13, 4.2) {half-equivalence:\\$\mathrm{pH} = \mathrm{p}K_a$};
  % equivalence point marker: V = 25 mL, steepest
  \addplot[engamber, only marks, mark=*, mark size=1.8pt]
    coordinates {(25, 7.5)};
  \draw[engamber, dotted, thick] (axis cs:25, 2) -- (axis cs:25, 7.5);
  \node[font=\footnotesize, align=left, text=engamber, anchor=south west]
    at (axis cs:25.5, 8.0) {equivalence point\\(steepest slope)};
  % buffer region bracket
  \draw[<->, engblack!60, thick] (axis cs:4, 3.4) -- (axis cs:21, 3.4);
  \node[font=\footnotesize, text=engblack!70, anchor=north]
    at (axis cs:12.5, 3.35) {buffer region (flat)};
\end{axis}
\end{tikzpicture}%
}%
\caption[An acid-base titration curve]{The pH of a weak acid as a
strong base is added, drawn against the volume of base. Because pH is
the negative base-ten logarithm of the hydrogen-ion concentration, the
vertical axis is already a logarithmic transform of the underlying
quantity, which is what turns a hundred-fold change in concentration
into a change of two pH units. The curve is sigmoidal: a flat buffer
region where the pH resists change and equals the acid's $\mathrm{p}K_a$
at the half-equivalence point, then a steep jump through the
equivalence point where a small further addition swings the pH by
several units. An engineer designing a buffer reads the target pH off
the flat region and the working capacity off its width; a control
engineer titrating to an endpoint reads the steep region as the
high-gain zone where the measurement is most sensitive and the control
hardest.}
\label{fig:vol02:ch02:titration-curve}
\end{figure}
